\chapter{Stratificazione, modelli di riferimento e incapsulamento}

Questo è forse il capitolo più importante di tutto il corso: introduce il \textbf{modello a
strati}, la chiave di lettura di tutto ciò che verrà dopo. Ogni parte successiva della dispensa si
occuperà di uno strato preciso.

% ══════════════════════════════════════════════════════════════
\section{Il Web è un'applicazione}
% ══════════════════════════════════════════════════════════════

Partiamo da un chiarimento: \textbf{il Web non è Internet}. Il Web è \emph{una} delle applicazioni
che girano sopra Internet; la posta elettronica, lo streaming, la messaggistica, il trasferimento di
file e i giochi sono altre applicazioni, e ciascuna chiede alla rete qualcosa di leggermente diverso.
Nel capitolo precedente abbiamo visto di cosa è fatta una rete e come i pezzi sono collegati. Ora
la domanda è:

\begin{center}
\large\itshape Che cosa deve esistere, e funzionare, perché una pagina compaia nel browser?
\end{center}

La risposta a questa domanda è la struttura di tutto il resto del corso.

\subsection{Cosa attraversa una singola richiesta HTTP}

Un browser sull'host A chiede una pagina al server Web B. Tra i due c'è un bel po' di roba:

\begin{figure}[H]
\centering
\begin{tikzpicture}[every node/.style={font=\footnotesize\sffamily}]
  \node (A) at (0,0) {\icon[1.2cm]{pc}};
  \node[below=-2pt of A] {A (browser)};
  \node (R) at (3.6,0) {\icon[1.4cm]{router}};
  \node[below=0pt of R] {router di casa};
  \node[box, fill=cloudbg, minimum width=1.4cm] (isp) at (8,0) {ISP};
  \node[netcloud, minimum width=3.2cm, minimum height=1.8cm] (inet) at (10.2,-2.5) {Internet};
  \node (B) at (13.3,-2.5) {\icon[0.8cm]{server}};
  \node[below=-2pt of B] {B (server Web)};
  \draw[wired] (A) -- node[above, lbl] {cavo Ethernet} (R);
  \draw[wired, netorange, line width=1.6pt] (R) -- node[above, lbl, text=netorange] {doppino telefonico} node[below, lbl, text=netorange] {e rete telefonica} (isp);
  \draw[wired, line width=2pt] (isp) -- (inet);
  \draw[wired, line width=2pt] (inet) -- (B);
  \draw[msg, netred, dashed] (0.3,1.0) .. controls (5,2.6) and (11,2.2) .. node[pos=0.45, above, lbl, text=netred] {richiesta HTTP di A verso B} (13,-1.6);
\end{tikzpicture}
\caption{Una richiesta HTTP attraversa mezzi diversi, apparati diversi e una maglia di reti di cui nessuno dei soggetti coinvolti conosce l'interno.}
\end{figure}

\begin{itemize}
    \item Ogni apparato sul cammino è \textbf{hardware più software}, costruito da un produttore diverso.
    \item Ogni tratto è un \textbf{mezzo diverso}, con i propri errori e la propria velocità.
    \item Nessuno può progettare, o anche solo descrivere, un sistema del genere \textbf{tutto in una volta}.
\end{itemize}

\dfn{Divide et impera}{
    Si spezza il sistema complesso in \textbf{sottosistemi}, ciascuno con un compito proprio e un
    modo ben definito di chiedere aiuto al successivo; poi si progetta un sottosistema alla volta.
    Nelle reti questa scomposizione ha un nome: \textbf{stratificazione} (\emph{layering}).
}

% ══════════════════════════════════════════════════════════════
\section{Internet vista “dai bulloni” e come servizio}
% ══════════════════════════════════════════════════════════════

Prima di entrare negli strati, ricapitoliamo i due modi complementari di guardare Internet.

\subsection{La vista “nuts and bolts”}

Dal punto di vista dei suoi componenti fisici (i ``dadi e bulloni''), Internet è fatta di:

\begin{itemize}
    \item \textbf{miliardi di dispositivi di calcolo} connessi, gli \textbf{host} o \emph{end system},
          che eseguono applicazioni di rete e stanno ai \textbf{bordi} (\emph{edge}) di Internet;
    \item \textbf{commutatori di pacchetto} (\emph{packet switch}), che inoltrano pacchetti (blocchi di
          dati): router e switch;
    \item \textbf{collegamenti di comunicazione}: fibra, rame, radio, satellite, ciascuno con il
          proprio \emph{transmission rate};
    \item \textbf{reti}: collezioni di dispositivi, router e collegamenti gestite da una singola organizzazione.
\end{itemize}

\begin{figure}[H]
    \centering
    \includegraphics[width=0.72\linewidth]{cap02/dispositivi_iot.png}
    \caption{Non solo computer: frigoriferi, cornici digitali, telecamere, pacemaker, materassi sensorizzati, orologi\dots{} sono tutti host di Internet.}
\end{figure}

In questa vista i \textbf{protocolli sono ovunque}: controllano l'invio e la ricezione dei messaggi
(HTTP per il Web, TCP, IP, Wi-Fi, 4G, Ethernet, i protocolli proprietari di Skype, \dots).

\subsection{Internet come infrastruttura di servizi}

\begin{figure}[H]
\centering
\begin{minipage}{0.42\linewidth}\centering
  \includegraphics[width=\linewidth]{cap02/internet_servizi.png}
\end{minipage}\hfill
\begin{minipage}{0.54\linewidth}
Vista dall'alto, Internet non è fatta di cavi: è un'\textbf{infrastruttura che fornisce servizi alle
applicazioni} (Web, video, posta, videoconferenza, commercio elettronico, social media, elettrodomestici
connessi).
\begin{itemize}
  \item Fornisce anche un'\textbf{interfaccia di programmazione}: i ``ganci'' (\emph{hooks}) con cui
        un'applicazione si aggancia a Internet per inviare e ricevere dati.
  \item Come il servizio postale (posta ordinaria, raccomandata, espresso), offre \textbf{opzioni di
        servizio}: veloce o economico, garantito o no.
\end{itemize}
Sono proprio quei ganci che programmeremo, in C, con i socket.
\end{minipage}
\caption{Protocolli e servizi in azione sulla rete di reti.}
\end{figure}

\info{L'applicazione vede solo un servizio}{
    Le applicazioni non devono preoccuparsi di come i dati attraversano la rete: il flusso dei dati
    viene offerto loro \textbf{come servizio} dallo strato inferiore (il trasporto). È esattamente
    questo il senso della stratificazione.
}

% ══════════════════════════════════════════════════════════════
\section{Protocolli e standard}
% ══════════════════════════════════════════════════════════════

Ricordiamo la definizione del capitolo precedente: un \textbf{protocollo} fissa il \textbf{formato}
e l'\textbf{ordine} dei messaggi, e le \textbf{azioni} da compiere quando un messaggio viene inviato o
ricevuto. Vale per le persone come per le macchine.

\begin{figure}[H]
    \centering
    \includegraphics[width=0.62\linewidth]{cap02/protocollo_umano_rete.png}
    \caption{Stessa forma: un saluto, una risposta, poi la richiesta vera e propria.}
\end{figure}

\subsection{Chi scrive gli standard}

\begin{itemize}
    \item L'\textbf{IETF} (\emph{Internet Engineering Task Force}) è una comunità aperta di esperti che
          definisce gli standard e i protocolli di Internet.
    \item Lavora in \textbf{working group aperti}, ciascuno dedicato a un aspetto di Internet.
    \item Ogni gruppo pubblica documenti numerati detti \textbf{RFC} (\emph{Request For Comments}), che
          descrivono e definiscono protocolli, formati e metodi. Per esempio IPv4 è definito
          nell'\textbf{RFC 791}, del 1981. Il nome è rimasto per ragioni storiche: nascevano come inviti
          alla discussione e sono diventati gli standard di Internet.
    \item L'\textbf{ISO} (\emph{International Standards Organization}) funziona in modo molto diverso: è
          un'organizzazione non basata su trattati internazionali, i cui membri sono gli enti nazionali
          di standardizzazione. È da qui che nasce il modello \textbf{OSI}.
\end{itemize}

\info{Come decide Internet}{
    Gli standard di Internet non nascono da un voto in commissione, ma dal principio
    \emph{rough consensus and running code}: consenso di massima e codice che funziona. Una specifica
    che nessuno implementa non diventa uno standard.
}

% ══════════════════════════════════════════════════════════════
\section{La stratificazione}
% ══════════════════════════════════════════════════════════════

Le reti sono complesse e fatte di molti pezzi: host, router, collegamenti su mezzi diversi,
applicazioni, protocolli, hardware, software. C'è speranza di organizzarne la struttura? Sì:
la risposta architetturale è una \textbf{pila di strati}, ciascuno costruito su quello sottostante.

\dfn{Modello a strati}{
    Un approccio architetturale in cui la comunicazione è realizzata da una pila di \textbf{strati}
    (\emph{layer}) tali che:
    \begin{itemize}
      \item ogni strato è un \textbf{livello di astrazione}: nasconde \emph{come} svolge il proprio compito;
      \item ogni strato \textbf{usa i servizi dello strato inferiore} e \textbf{fornisce servizi allo
            strato superiore};
      \item tra due strati adiacenti c'è un'\textbf{interfaccia}: l'insieme delle operazioni che lo
            strato inferiore mette a disposizione di quello superiore.
    \end{itemize}
    La cima della pila è rivolta alle applicazioni (alto livello), il fondo all'hardware (basso livello).
}

\begin{figure}[H]
\centering
\begin{minipage}{0.33\linewidth}\centering
  \includegraphics[width=0.85\linewidth]{cap02/stratificazione.png}
\end{minipage}
\begin{minipage}{0.55\linewidth}\centering
\begin{tikzpicture}[every node/.style={font=\small\sffamily}]
  \node[lay app, minimum width=3.4cm] (n1) at (0,1.4) {strato $n+1$};
  \node[lay tra, minimum width=3.4cm] (n) at (0,0) {strato $n$};
  \node[lay net, minimum width=3.4cm] (n0) at (0,-1.4) {strato $n-1$};
  \draw[{Stealth}-{Stealth}, thick] (n1) -- node[right, lbl] {interfaccia $n/n+1$} (n);
  \draw[{Stealth}-{Stealth}, thick] (n) -- node[right, lbl] {interfaccia $n-1/n$} (n0);
  \node[lbl, text width=3cm, align=left] at (-4,0.7) {lo strato $n$ \textbf{fornisce} servizi allo strato $n+1$\dots};
  \node[lbl, text width=3cm, align=left] at (-4,-0.7) {\dots\ \textbf{usando} i servizi dello strato $n-1$};
  \draw[-{Stealth}, primary, thick] (2.6,2) -- node[right, lbl] {verso le applicazioni} (2.6,2.6);
  \draw[-{Stealth}, primary, thick] (2.6,-2) -- node[right, lbl] {verso l'hardware} (2.6,-2.6);
\end{tikzpicture}
\end{minipage}
\caption{Ogni strato si appoggia a quello sotto e serve quello sopra, attraverso un'interfaccia.}
\end{figure}

\subsection{Moduli, interfacce e incapsulamento delle funzioni}

L'idea non è nuova: è la stessa che si usa per progettare qualunque software di una certa dimensione. Quando
un sistema viene diviso in sottosistemi, per ciascuno si stabilisce un'\textbf{interfaccia}, cioè l'insieme
dei servizi che il modulo offre all'esterno; \emph{come} quei servizi siano realizzati resta nascosto dentro
il modulo. È l'\textbf{incapsulamento delle funzioni}: ogni modulo può chiamare i servizi di un altro per il
proprio funzionamento, senza sapere nulla della sua implementazione. Chi ha scritto classi in C++ o in Java
ha già fatto esattamente questo: metodi pubblici all'esterno, dettagli privati all'interno.

\begin{figure}[H]
\centering
\begin{tikzpicture}[every node/.style={font=\footnotesize\sffamily}]
  % modulo A (offre servizi)
  \node[draw=netblue, thick, minimum width=2.8cm, minimum height=1.6cm, fill=netblue!5] (A) at (0,0) {};
  \node[text=netblue!70!black, font=\small\bfseries\sffamily] at (0,0.45) {Modulo A};
  \node[draw=black!30, dashed, fill=black!5, minimum width=2.2cm, minimum height=0.6cm, font=\scriptsize\sffamily] at (0,-0.25) {implementazione};
  \foreach \y/\n [count=\i] in {0.5/apri, 0/scrivi, -0.5/chiudi} {
    \draw[thick, netblue] (1.4,\y) -- (1.85,\y) node[pinint, draw=netblue] (pa\i) {};
    \node[anchor=west, text=netblue!70!black] at (1.95,\y) {\texttt{\n()}};
  }
  % modulo B (usa i servizi)
  \node[draw=netgreen!70!black, thick, minimum width=2.8cm, minimum height=1.6cm, fill=netgreen!5] (B) at (8,0) {};
  \node[text=netgreen!50!black, font=\small\bfseries\sffamily] at (8,0.45) {Modulo B};
  \node[draw=black!30, dashed, fill=black!5, minimum width=2.2cm, minimum height=0.6cm, font=\scriptsize\sffamily] at (8,-0.25) {implementazione};
  \draw[-{Stealth}, thick, netgreen!60!black] ([yshift=0.1cm]B.west) -- node[above, lbl] {chiama \texttt{scrivi()}} (3.5,0);
  \draw[decorate, decoration={brace, amplitude=4pt}] (1.85,-0.85) -- (1.85,0.85);
  \node[lbl, text=netblue!70!black, anchor=north] at (2.4,-0.95) {interfaccia: i ``piedini'' del modulo};
  \node[lbl, text width=4.2cm, anchor=north] at (8,-1.0) {B conosce solo i piedini di A, non ciò che c'è dentro};
\end{tikzpicture}
\caption{Ogni modulo espone all'esterno i propri servizi (i ``piedini'') e nasconde l'implementazione: è il disegno fatto alla lavagna per introdurre gli strati.}
\end{figure}

Nell'architettura a strati questa idea diventa più \textbf{rigida}: i moduli sono impilati, e ogni strato
offre i propri servizi allo strato \textbf{immediatamente superiore} e usa quelli dello strato
\textbf{immediatamente inferiore} per implementarli. Lo strato in cima è quello delle applicazioni (o offerto
alle applicazioni); quello più in basso è l'hardware, lo strato fisico.

\info{Architetture chiuse e aperte}{
    Nella concezione di base l'architettura a strati è \textbf{chiusa}: uno strato può rivolgersi soltanto a
    quello subito sotto. In alcuni casi, per efficienza o per necessità, la si \textbf{apre}, permettendo a uno
    strato di usare direttamente servizi di strati non adiacenti. È un'eccezione, e ha un prezzo: i moduli
    diventano meno indipendenti.
}

\subsection{La stratificazione non è un'idea delle reti}

Pile di strati le usate già, ovunque un sistema sia diventato troppo grande per essere progettato
in un pezzo solo.

\begin{figure}[H]
\centering
\begin{tikzpicture}[every node/.style={font=\small\sffamily}]
  \begin{scope}
    \node[font=\bfseries\sffamily] at (0,3.1) {Un sistema operativo};
    \node[layer, fill=lapp!60, minimum width=3.6cm] (a) at (0,2.2) {Applicazioni};
    \node[layer, fill=ltra!60, minimum width=3.6cm] (o) at (0,1.2) {Sistema operativo};
    \node[layer, fill=lnet!60, minimum width=3.6cm] (d) at (0,0.2) {Driver};
    \node[layer, fill=lphy, minimum width=3.6cm] (h) at (0,-0.8) {Hardware};
    \node[font=\ttfamily\small] (vim) at (3.2,2.9) {vim};
    \draw[-{Stealth}, netred, thick] (vim) to[out=-90, in=0] node[right, lbl, text=netred] {\texttt{write}} (a.east);
    \draw[-{Stealth}, netred, thick] ([xshift=1.2cm]a.south) -- ([xshift=1.2cm]o.north);
    \draw[-{Stealth}, netred, thick] ([xshift=1.2cm]o.south) -- ([xshift=1.2cm]d.north);
    \draw[-{Stealth}, netred, thick] ([xshift=1.2cm]d.south) -- ([xshift=1.2cm]h.north);
    \node[lbl, text width=4cm] at (0,-1.7) {una sola \texttt{write} attraversa quattro interfacce};
  \end{scope}
  \begin{scope}[xshift=8.5cm]
    \node[font=\bfseries\sffamily] at (0,3.1) {Un sistema informativo};
    \node[layer, fill=lapp!60, minimum width=3.6cm] (ui) at (0,2.2) {Interfaccia utente};
    \node[layer, fill=ltra!60, minimum width=3.6cm] (la) at (0,1.2) {Logica applicativa};
    \node[layer, fill=lnet!60, minimum width=3.6cm] (gd) at (0,0.2) {Gestione dei dati};
    \node[cylinder, shape border rotate=90, aspect=0.25, draw, fill=lphy, minimum height=0.9cm, minimum width=0.8cm] (db) at (2.8,0.2) {DB};
    \draw[thick] (gd) -- (db);
    \node at (-2.8,2.8) {\icon[0.7cm]{laptop}};
    \draw[thick, -{Stealth}] (-2.6,2.5) -- (ui.west);
    \node[lbl, text width=4.4cm] at (0,-0.9) {architettura a tre strati: l'utente non tocca mai direttamente il database};
  \end{scope}
\end{tikzpicture}
\caption{Due pile di strati che conoscete già.}
\end{figure}

Le reti sono semplicemente il caso più grande di tutti.

\subsection{Cosa guadagniamo, e cosa paghiamo}

\begin{itemize}
    \item \textbf{Struttura}: una struttura esplicita permette di \textbf{dare un nome} ai pezzi del
          sistema e alle loro relazioni; è un \emph{modello di riferimento} di cui discutere.
    \item \textbf{Modularità}: l'implementazione di uno strato può cambiare senza che nient'altro si
          muova, \textbf{purché il servizio offerto resti lo stesso}. Sostituire tutte le linee
          telefoniche con canali satellitari è una modifica interna a un solo strato: gli strati
          superiori non se ne accorgono.
\end{itemize}

\begin{esempio}{La modularità nella vita di tutti i giorni}
\begin{itemize}
    \item Cambiate stampante, e con lei il \textbf{driver}: il sistema operativo continua a mandare le stampe
          esattamente come prima, e le applicazioni non si accorgono di nulla.
    \item State scrivendo un messaggio collegati al Wi-Fi di casa, poi uscite e il telefono passa alla
          \textbf{rete cellulare}. L'applicazione di messaggistica si rivolge sempre al sistema operativo nello
          stesso modo: del cambio di tecnologia si occupano gli strati inferiori.
\end{itemize}
\end{esempio}

\info{L'analogia del viaggio aereo}{
    Un viaggio in aereo è una pila di servizi: biglietto (acquisto / reclamo), bagagli (imbarco /
    ritiro), gate (imbarco / sbarco), pista (decollo / atterraggio), rotta dell'aereo. Ogni ``strato''
    dell'aeroporto di partenza ha il suo corrispondente in quello di arrivo, e se cambia la procedura
    ai gate il resto del sistema non se ne accorge.
}

\warning{Il prezzo dell'architettura a strati}{
    La domanda ``\emph{layering considered harmful?}'' non è retorica. Alla lavagna il prezzo è stato riassunto
    in tre voci:
    \begin{itemize}
      \item \textbf{maggiore complessità}: tutte le \textbf{interfacce} tra gli strati vanno progettate e
            implementate, e ogni strato aggiunge le proprie informazioni di controllo (\textbf{overhead});
      \item \textbf{eventuali ridondanze}: uno strato può rifare qualcosa che uno strato inferiore fa già, per
            esempio la rilevazione e la correzione degli errori, svolta sia su ogni singolo collegamento sia da
            un estremo all'altro. Più calcolo, a volte inutile;
      \item \textbf{rigidità}: le informazioni interne di ogni strato sono \textbf{nascoste}. Quando poi uno
            strato ha bisogno di un'informazione che sta in un altro (succede spesso nelle reti), le strade sono
            due: \textbf{riprogettare} gli strati, oppure adottare soluzioni ``\emph{quick and dirty}'', rapide,
            efficaci e un po' sporche, come vedremo per esempio con l'uso delle porte.
    \end{itemize}
    È un compromesso: si accetta un po' di inefficienza in cambio di un'enorme semplicità di progetto.
}

% ══════════════════════════════════════════════════════════════
\section{Servizio, interfaccia e protocollo}
% ══════════════════════════════════════════════════════════════

Tre concetti da tenere rigorosamente separati.

\dfn{Servizio e protocollo}{
    Un \textbf{servizio} è l'insieme delle operazioni che uno strato offre allo strato superiore: dice
    \textbf{che cosa} fa lo strato, mai \textbf{come} lo fa.

    Un \textbf{protocollo} è l'insieme delle regole sul formato e sul significato dei messaggi
    scambiati con l'\textbf{entità paritaria} (\emph{peer entity}), cioè lo stesso strato sull'altra
    macchina.
}

\begin{figure}[H]
    \centering
    \includegraphics[width=0.62\linewidth]{cap02/servizio_protocollo.png}
    \caption{I servizi stanno sulle interfacce verticali, i protocolli sullo scambio orizzontale.}
\end{figure}

L'\textbf{interfaccia} è il punto in cui il servizio viene offerto: l'insieme delle operazioni che
lo strato inferiore rende disponibili a quello superiore. In sintesi:

\begin{center}
\begin{tabular}{lll}
\toprule
\textbf{Concetto} & \textbf{Direzione} & \textbf{Risponde alla domanda} \\
\midrule
Servizio & verticale (tra strati adiacenti) & \emph{che cosa} offre lo strato? \\
Interfaccia & verticale & \emph{come si chiede} il servizio? \\
Protocollo & orizzontale (tra strati omologhi) & \emph{come} lo strato realizza il servizio? \\
\bottomrule
\end{tabular}
\end{center}

\dfn{Entità e protocolli}{
    Le \textbf{entità} (gli elementi attivi di uno strato, hardware o software) \textbf{usano i
    protocolli per implementare i servizi} che offrono.
}

\warning{Servizio e protocollo sono indipendenti}{
    Uno strato può cambiare completamente protocollo senza che gli strati superiori se ne accorgano,
    purché il servizio offerto resti lo stesso. Confondere servizio e protocollo è un errore tipico.
}

\begin{esempio}{Un messaggio alla zia}
Premete ``invia'' sulla vostra applicazione di messaggistica. L'applicazione non trasmette nulla da sé: chiama
il servizio di invio che il \textbf{livello di trasporto} del telefono offre attraverso la propria interfaccia,
e da quel momento se ne disinteressa. Il livello di trasporto dialoga, con il \textbf{protocollo di trasporto},
con l'entità di trasporto sul telefono della zia; per farlo usa i servizi del \textbf{livello di rete}, che a
sua volta dialoga con il livello di rete dall'altra parte, occupandosi di attraversare la topologia di Internet.
\begin{center}
\begin{tikzpicture}[every node/.style={font=\footnotesize\sffamily}]
  % mittente
  \pic[black!80] at (-2.6,1.55) {persona};
  \node[anchor=east] at (-2.9,1.6) {voi};
  \node[lay app, minimum width=2.3cm] (a1) at (0,1.8) {Applicazione};
  \node[lay tra, minimum width=2.3cm] (t1) at (0,0.3) {Trasporto};
  \node[lay net, minimum width=2.3cm] (n1) at (0,-1.2) {Rete};
  \draw[-{Stealth}, thick] (-2.2,1.8) -- node[above, lbl] {``invia''} (a1.west);
  % destinataria
  \pic[black!80] at (11.9,1.55) {personacapelli};
  \node[anchor=west] at (12.2,1.6) {la zia};
  \node[lay app, minimum width=2.3cm] (a2) at (9.3,1.8) {Applicazione};
  \node[lay tra, minimum width=2.3cm] (t2) at (9.3,0.3) {Trasporto};
  \node[lay net, minimum width=2.3cm] (n2) at (9.3,-1.2) {Rete};
  \draw[-{Stealth}, thick] (a2.east) -- node[above, lbl] {legge} (11.5,1.8);
  % piedini delle interfacce
  \foreach \st in {t1,n1,t2,n2} \foreach \dx in {-0.4,0,0.4} \draw[thick] ([xshift=\dx cm]\st.north) -- ++(0,0.22) node[pinint] {};
  \draw[-{Stealth}, very thick, netred] ([xshift=0.8cm]a1.south) -- node[right, lbl, text=netred] {invio messaggio} ([xshift=0.8cm]t1.north);
  \draw[-{Stealth}, very thick, netred] ([xshift=0.8cm]t1.south) -- ([xshift=0.8cm]n1.north);
  \draw[-{Stealth}, very thick, netgreen!60!black] ([xshift=-0.8cm]n2.north) -- ([xshift=-0.8cm]t2.south);
  \draw[-{Stealth}, very thick, netgreen!60!black] ([xshift=-0.8cm]t2.north) -- ([xshift=-0.8cm]a2.south);
  % protocolli
  \node[cloud, cloud puffs=12, aspect=3.2, draw=netorange, fill=netorange!10, thick, minimum width=2.6cm, minimum height=0.8cm, inner sep=1pt] (pr) at (4.65,0.3) {protocollo di trasporto};
  \draw[{Stealth}-{Stealth}, dashed, thick, netblue] (t1.east) -- (pr.west);
  \draw[{Stealth}-{Stealth}, dashed, thick, netblue] (pr.east) -- (t2.west);
  \draw[{Stealth}-{Stealth}, dashed, thick, netgreen!60!black] (n1.east) -- node[above, lbl] {protocollo di rete (attraverso la topologia di Internet)} (n2.west);
  \node[lbl, text=black!60, anchor=north] at (0,-1.65) {interfaccia = servizi dello strato sotto};
\end{tikzpicture}
\end{center}
Tra uno strato e quello sotto, \textbf{nello stesso telefono}, c'è un'\textbf{interfaccia} (i piedini); tra due
entità dello stesso strato, \textbf{su telefoni diversi}, c'è un \textbf{protocollo} (le frecce tratteggiate).
\end{esempio}

\subsection{Strati omologhi: comunicazione virtuale e fisica}

\begin{figure}[H]
    \centering
    \includegraphics[width=0.62\linewidth]{cap02/strati_omologhi.png}
    \caption{Lo strato $n$ del mittente dialoga logicamente con lo strato $n$ del destinatario.}
\end{figure}

\dfn{Comunicazione orizzontale}{
    Quando due dispositivi comunicano, \textbf{ogni strato del mittente comunica con lo stesso strato
    del destinatario}, e con nient'altro, per mezzo di uno specifico protocollo.
}

\begin{itemize}
    \item Quella freccia orizzontale \textbf{non è un cavo}: è una conversazione \emph{logica}.
    \item I dati \textbf{scendono} realmente lungo la pila del mittente e \textbf{risalgono} quella del
          destinatario. Solo al livello fisico qualcosa viaggia davvero.
    \item La catena si ripete a ogni strato: per fornire i propri servizi uno strato deve comunicare con il suo
          pari, e per farlo chiama i servizi dello strato inferiore, che a sua volta comunica con il proprio pari
          usando un protocollo \textbf{diverso}, e così via fino al protocollo fisico, che manda segnali su un mezzo.
\end{itemize}

\warning{Tutte queste trasmissioni sono virtuali}{
    Nelle reti c'è comunicazione tra due macchine a \textbf{ogni} strato, con i relativi protocolli. Ma
    tutte queste trasmissioni sono \textbf{virtuali}, salvo quella a livello fisico, che è l'unica a
    comportare effettivi fenomeni fisici. Il livello di trasporto del mittente ``crede'' di parlare
    direttamente con il trasporto del destinatario, ma i suoi dati scendono lungo tutta la pila,
    viaggiano sul filo e risalgono dall'altra parte.
}

\begin{esempio}{L'analogia postale}
Il modello a strati funziona come il servizio postale: ogni strato fa un solo lavoro, si affida allo strato
sotto e parla il proprio ``linguaggio''; ogni indirizzo è letto dallo strato per cui è stato scritto, e da
nessun altro.
\begin{center}
\begin{tikzpicture}[every node/.style={font=\footnotesize\sffamily},
    pbox/.style={draw=black!60, rounded corners=2pt, minimum width=3.3cm, minimum height=0.95cm, align=center}]
  % colonna sinistra
  \node[pbox, fill=lapp!60] (m) at (0,0) {\textbf{Mittente}\\scrive la lettera};
  \node[pbox, fill=ltra!60] (u1) at (0,-1.9) {\textbf{Ufficio postale}\\raccolta e smistamento};
  \node[pbox, fill=lnet!60] (x1) at (0,-3.8) {\textbf{Centro di trasporto}\\treni, aerei, navi};
  % colonna destra
  \node[pbox, fill=lapp!60] (d) at (11,0) {\textbf{Destinatario}\\legge la lettera};
  \node[pbox, fill=ltra!60] (u2) at (11,-1.9) {\textbf{Ufficio postale}\\postino (bici, motorino, auto)};
  \node[pbox, fill=lnet!60] (x2) at (11,-3.8) {\textbf{Centro di trasporto}\\treni, aerei, navi};
  % interfacce verticali
  \draw[-{Stealth}, thick, netred] (m) -- node[left, lbl, text=netred] {buca delle lettere} (u1);
  \draw[-{Stealth}, thick, netred] (u1) -- node[left, lbl, text=netred] {spedizione} (x1);
  \draw[-{Stealth}, thick, netgreen!60!black] (x2) -- node[right, lbl, text=netgreen!50!black] {arrivo} (u2);
  \draw[-{Stealth}, thick, netgreen!60!black] (u2) -- node[right, lbl, text=netgreen!50!black] {cassetta del destinatario} (d);
  % dialoghi orizzontali
  \draw[{Stealth}-{Stealth}, dashed, thick] (m) -- node[above, lbl] {scrivere e leggere una lettera} (d);
  \draw[{Stealth}-{Stealth}, dashed, thick] (u1) -- node[above, lbl] {usano l'\textbf{indirizzo} di mittente e destinatario} (u2);
  \draw[very thick, black!70] (x1) -- node[above, lbl, text width=6.5cm] {trasporto fisico, che ragiona sugli indirizzi dei \textbf{due uffici postali}} (x2);
\end{tikzpicture}
\end{center}
\begin{itemize}
    \item \textbf{Strato superiore}: mittente e destinatario vedono soltanto due buche per le lettere, una in cui
          infilare la lettera e una da cui ritirarla. I loro ``protocolli'' sono scrivere e leggere una lettera.
    \item \textbf{Strato intermedio}: gli uffici postali raccolgono la lettera e, dall'altra parte, la fanno
          consegnare dal postino. Il loro linguaggio sono gli \textbf{indirizzi} del mittente e del destinatario.
    \item \textbf{Strato inferiore}: treni, aerei e navi portano fisicamente la posta da un ufficio all'altro, e
          ragionano solo in termini di uffici di provenienza e di destinazione.
\end{itemize}
Tanenbaum propone un'analogia altrettanto efficace: un filosofo che detta una lettera a un traduttore, il
quale la passa a una segretaria che la spedisce; dall'altra parte la catena si ripercorre al contrario, e ogni
figura ``parla'' solo con la figura corrispondente.
\end{esempio}

% ══════════════════════════════════════════════════════════════
\section{Servizi con e senza connessione}
% ══════════════════════════════════════════════════════════════

Prima di vedere gli strati concreti, serve una distinzione \textbf{ortogonale} a quella degli strati:
quella tra servizi \textbf{orientati alla connessione} e servizi \textbf{senza connessione}.

\subsection{Servizio orientato alla connessione}

\dfn{Servizio orientato alla connessione}{
    Modellato sul \textbf{sistema telefonico}: l'utente del servizio prima \textbf{stabilisce} una
    connessione, poi la \textbf{usa}, infine la \textbf{rilascia}.
}

\begin{itemize}
    \item La connessione si comporta come un \textbf{tubo}: ciò che entra da un'estremità esce
          dall'altra, normalmente \textbf{nello stesso ordine}.
    \item L'apertura permette una \textbf{negoziazione} dei parametri (dimensione massima dei messaggi,
          qualità del servizio, \dots): una parte propone, l'altra accetta o fa una controproposta.
    \item Una connessione che riserva anche delle \textbf{risorse}, per esempio una velocità fissa, si
          chiama \textbf{circuito}: un cammino sulla rete su cui viaggiano tutti i dati della sessione.
    \item Conviene quando c'è una quantità significativa di dati da scambiare.
\end{itemize}

E se si ha un solo messaggio breve da inviare, vale la pena aprire e chiudere una connessione?

\subsection{Servizio senza connessione}

\dfn{Servizio senza connessione (connectionless)}{
    Modellato sul \textbf{sistema postale}: non si stabilisce nulla prima di trasmettere. Ogni
    messaggio porta con sé l'\textbf{indirizzo completo} di destinazione e viene \textbf{gestito
    individualmente} dai nodi intermedi.
}

\begin{figure}[H]
\centering
\begin{tikzpicture}[every node/.style={font=\footnotesize\sffamily}]
  \node (A) at (0,0) {\icon[0.9cm]{pc}}; \node[below=-2pt of A] {A};
  \node (B) at (10,0) {\icon[1cm]{laptop}}; \node[below=-2pt of B] {B};
  \node[cloud, cloud puffs=14, aspect=2.6, draw=netblue!50, fill=cloudbg, minimum width=6.2cm, minimum height=3cm] at (5,0) {};
  \node[gnode] (n1) at (3,0.5) {}; \node[gnode] (n2) at (4.5,1) {}; \node[gnode] (n3) at (4.3,-0.8) {};
  \node[gnode] (n4) at (6,0.1) {}; \node[gnode] (n5) at (7,0.9) {}; \node[gnode] (n6) at (6.9,-0.8) {};
  \foreach \a/\b in {n1/n2, n1/n3, n2/n4, n3/n4, n2/n5, n4/n5, n4/n6, n3/n6, n5/n6} \draw[black!40, thick] (\a) -- (\b);
  \draw[wired] (A) -- (n1); \draw[wired] (n5) -- (B); \draw[wired] (n6) -- (B);
  \draw[netred, very thick, -{Stealth}] (n1) -- (n2) -- (n5);
  \draw[netgreen, very thick, -{Stealth}] (n1) -- (n3) -- (n6);
  \node[text=netred] at (3.4,1.55) {pacchetto 1};
  \node[text=netgreen] at (5.3,-1.45) {pacchetto 2};
  \fill[netred] (1.2,0.35) rectangle (1.6,0.6); \fill[netgreen] (1.8,0.35) rectangle (2.2,0.6);
  \node at (5,-2.1) {ogni pacchetto è gestito individualmente: può seguire un cammino diverso e arrivare fuori ordine};
\end{tikzpicture}
\caption{In un servizio senza connessione i pacchetti tra gli stessi due estremi possono seguire strade diverse.}
\end{figure}

\begin{itemize}
    \item Tra gli stessi due estremi la rete offre più cammini, quindi due messaggi possono seguirne
          di diversi e \textbf{arrivare fuori ordine}.
    \item Nella commutazione \textbf{store-and-forward} un nodo riceve un messaggio \textbf{per intero}
          prima di inoltrarlo al nodo successivo.
    \item Nella commutazione \textbf{cut-through} un nodo comincia a inoltrare il messaggio \textbf{prima}
          di averlo ricevuto completamente (appena ha letto l'indirizzo di destinazione).
\end{itemize}

\begin{figure}[H]
\centering
\begin{tikzpicture}[every node/.style={font=\footnotesize\sffamily}]
  \begin{scope}
    \node[font=\bfseries\sffamily] at (2.5,0.7) {Store-and-forward};
    \foreach \x/\n in {0/in, 2.5/nodo, 5/out} { \draw[thick] (\x,0) -- (\x,-4); }
    \node at (0,0.25) {ingresso}; \node at (2.5,0.25) {nodo}; \node at (5,0.25) {uscita};
    \fill[netblue!40] (0,0) -- (2.5,-0.4) -- (2.5,-1.6) -- (0,-1.2) -- cycle;
    \fill[netblue!40] (2.5,-1.6) -- (5,-2.0) -- (5,-3.2) -- (2.5,-2.8) -- cycle;
    \node[lbl, text width=2cm, fill=white, inner sep=1pt] at (3.8,-0.8) {aspetta l'ultimo bit};
  \end{scope}
  \begin{scope}[xshift=8cm]
    \node[font=\bfseries\sffamily] at (2.5,0.7) {Cut-through};
    \foreach \x in {0,2.5,5} { \draw[thick] (\x,0) -- (\x,-4); }
    \node at (0,0.25) {ingresso}; \node at (2.5,0.25) {nodo}; \node at (5,0.25) {uscita};
    \fill[netblue!40] (0,0) -- (2.5,-0.4) -- (2.5,-1.6) -- (0,-1.2) -- cycle;
    \fill[netblue!40] (2.5,-0.6) -- (5,-1.0) -- (5,-2.2) -- (2.5,-1.8) -- cycle;
    \draw[netred, thick] (2.35,-0.6) -- (2.65,-0.6);
    \node[lbl, text width=2.3cm, fill=white, inner sep=1pt, text=netred] at (3.9,-2.9) {inizia a inoltrare appena letta l'intestazione};
  \end{scope}
\end{tikzpicture}
\caption{Store-and-forward contro cut-through: il secondo riduce il ritardo, ma può inoltrare anche messaggi che si riveleranno danneggiati.}
\end{figure}

\dfn{Datagram}{
    Un servizio \textbf{senza connessione e inaffidabile}, cioè senza conferma di ricezione
    (\emph{unacknowledged}). Il nome richiama il \emph{telegramma}, che anch'esso non prevedeva
    ricevuta di ritorno.
}

\subsection{Affidabilità}

\dfn{Affidabilità}{
    Un servizio è \textbf{affidabile} se il ricevente \textbf{conferma} (\emph{acknowledgement}, ACK)
    ciò che ha ricevuto, così che il mittente sappia che è arrivato. L'affidabilità è
    \textbf{ortogonale} alla connessione: sia i servizi con connessione sia quelli senza possono essere
    affidabili o no.
}

Perché mai si dovrebbe rinunciare all'affidabilità?

\begin{itemize}
    \item Le conferme costano \textbf{overhead e ritardo}, e a volte il prezzo è troppo alto.
    \item L'affidabilità può semplicemente \textbf{non essere disponibile}: Ethernet non la fornisce, e
          sono gli strati superiori a dover recuperare gli errori.
    \item Nel \textbf{Voice over IP} confermare ogni pacchetto e ripristinarne l'ordine sarebbe non solo inutile ma
          dannoso: un ritardo di tre o quattro secondi rende una telefonata inutilizzabile, mentre una breve
          interruzione dell'audio si rimedia ripetendo una frase. Tutte le applicazioni VoIP cercano di
          \textbf{ridurre il ritardo} e accettano i buchi: l'inaffidabilità è la scelta \emph{giusta}, non un
          compromesso. I parametri della comunicazione vanno adattati alle esigenze dell'applicazione.
\end{itemize}

\info{Affidabile sopra inaffidabile}{
    La maggior parte dei servizi affidabili che vedremo nel corso è costruita \textbf{sopra} un servizio
    inaffidabile: TCP (affidabile) sopra IP (inaffidabile) ne è l'esempio principale.
}

\subsection{Sei tipi di servizio}

\begin{table}[H]
\centering
\begin{tabular}{lll}
\toprule
\textbf{Tipo} & \textbf{Servizio} & \textbf{Esempio} \\
\midrule
\multirow{3}{*}{Orientato alla connessione}
  & Flusso di messaggi affidabile & Sequenza di pagine \\
  & Flusso di byte affidabile & Download di un film \\
  & Connessione inaffidabile & Voice over IP \\
\midrule
\multirow{3}{*}{Senza connessione}
  & Datagram inaffidabile & Posta elettronica ``spazzatura'' \\
  & Datagram con conferma & Messaggi con le ``spunte'' di consegna \\
  & Richiesta-risposta (\emph{request-reply}) & Interrogazione a un database \\
\bottomrule
\end{tabular}
\caption{Sei servizi comuni, con e senza connessione.}
\end{table}

\info{Le due righe da ricordare}{
    Il \textbf{flusso di byte affidabile} è il servizio che ci darà \textbf{TCP}; il \textbf{datagram
    inaffidabile} è quello che ci darà \textbf{UDP}.
}

\subsection{Le primitive di servizio}

Un servizio si specifica formalmente tramite le \textbf{primitive}, cioè le operazioni con cui un
processo lo richiede. Quando la pila protocollare è implementata nel sistema operativo, le primitive
sono \textbf{chiamate di sistema} (\emph{system call}).

\begin{table}[H]
\centering
\begin{tabular}{ll}
\toprule
\textbf{Primitiva} & \textbf{Significato} \\
\midrule
\texttt{LISTEN} & si blocca in attesa di una connessione in ingresso \\
\texttt{CONNECT} & stabilisce una connessione con un peer in attesa \\
\texttt{ACCEPT} & accetta una connessione in ingresso da un peer \\
\texttt{RECEIVE} & si blocca in attesa di un messaggio in ingresso \\
\texttt{SEND} & invia un messaggio al peer \\
\texttt{DISCONNECT} & termina la connessione \\
\bottomrule
\end{tabular}
\caption{Le primitive di un semplice servizio orientato alla connessione.}
\end{table}

\begin{figure}[H]
\centering
\begin{tikzpicture}[yscale=0.8, every node/.style={font=\footnotesize\sffamily}]
  \timeline{c}{0}{Client}{7.3}
  \timeline{s}{8}{Server}{7.3}
  \node[right, font=\ttfamily\footnotesize] at (8,-0.35) {LISTEN};
  \draw[msg] (0,-0.8) node[left, font=\ttfamily\footnotesize] {CONNECT} -- node[above, sloped] {1. richiesta di connessione} (8,-1.6);
  \draw[msg] (8,-2.0) node[right, font=\ttfamily\footnotesize] {ACCEPT} -- node[above, sloped] {2. connessione accettata} (0,-2.8);
  \draw[msg, netblue] (0,-3.2) node[left, font=\ttfamily\footnotesize] {SEND} -- node[above, sloped] {3. richiesta di dati} (8,-4.0) node[right, font=\ttfamily\footnotesize] {RECEIVE};
  \draw[msg, netgreen] (8,-4.4) node[right, font=\ttfamily\footnotesize] {SEND} -- node[above, sloped] {4. risposta} (0,-5.2) node[left, font=\ttfamily\footnotesize] {RECEIVE};
  \draw[msg, netred] (0,-5.6) node[left, font=\ttfamily\footnotesize] {DISCONNECT} -- node[above, sloped] {5. rilascio} (8,-6.4);
  \draw[msg, netred] (8,-6.6) node[right, font=\ttfamily\footnotesize] {DISCONNECT} -- node[above, sloped] {6. rilascio confermato} (0,-7.2);
\end{tikzpicture}
\caption{Un'interazione client-server realizzata con le sei primitive.}
\end{figure}

\info{Le ritroveremo in C}{
    Guardatele bene: sono esattamente l'interfaccia su cui scriveremo i programmi in C nel capitolo
    sui socket: \texttt{listen()}, \texttt{connect()}, \texttt{accept()}, \texttt{send()},
    \texttt{read()}, \texttt{close()}.
}

% ══════════════════════════════════════════════════════════════
\section{Il modello di riferimento ISO/OSI}
% ══════════════════════════════════════════════════════════════

Negli anni '80 l'\textbf{ISO} definì un modello a strati di riferimento per le reti di calcolatori: il
modello \textbf{OSI} (\emph{Open Systems Interconnection}), cioè sistemi \emph{aperti} alla
comunicazione con altri sistemi.

\begin{figure}[H]
\centering
\begin{minipage}{0.42\linewidth}\centering
  \includegraphics[width=\linewidth]{cap02/osi_sette_strati.png}
\end{minipage}\hfill
\begin{minipage}{0.54\linewidth}
\begin{itemize}
  \item \textbf{Sette strati}, ciascuno creato dove serviva un'\textbf{astrazione diversa}.
  \item Il suo contributo più importante è aver reso \textbf{servizi}, \textbf{interfacce} e
        \textbf{protocolli} tre concetti \textbf{esplicitamente separati}.
  \item Il mittente scende lungo la pila, il destinatario la risale.
\end{itemize}
\end{minipage}
\caption{I sette strati del modello ISO/OSI.}
\end{figure}

\begin{table}[H]
\centering
\begin{tabular}{clp{9.5cm}}
\toprule
\textbf{\#} & \textbf{Strato} & \textbf{Compito} \\
\midrule
7 & Applicazione & servizi all'utente: login remoto, trasferimento file, posta, Web \\
6 & Presentazione & rappresentazione dei dati: traduzione tra codifiche, cifratura, compressione \\
5 & Sessione & controllo del dialogo (chi parla, quando e per quanto) e checkpoint \\
4 & Trasporto & consegna end-to-end al processo giusto (porte), segmentazione, controllo di flusso e degli errori \\
3 & Rete & consegna dei pacchetti dall'host sorgente all'host destinazione: indirizzamento logico e routing \\
2 & Collegamento & consegna dei frame da un nodo al successivo (un hop), controllo degli errori e accesso al mezzo \\
1 & Fisico & trasmissione dei singoli bit sul mezzo \\
\bottomrule
\end{tabular}
\caption{I sette strati OSI, dall'alto verso il basso.}
\end{table}

\info{Un trucco mnemonico}{
    Dal basso verso l'alto: \textbf{F}isico, \textbf{C}ollegamento, \textbf{R}ete, \textbf{T}rasporto,
    \textbf{S}essione, \textbf{P}resentazione, \textbf{A}pplicazione. In inglese si usa la frase
    ``\emph{\textbf{P}lease \textbf{D}o \textbf{N}ot \textbf{T}hrow \textbf{S}ausage \textbf{P}izza
    \textbf{A}way}'' (Physical, Data link, Network, Transport, Session, Presentation, Application).
}

\subsection{Non tutti i nodi implementano tutti gli strati}

\begin{figure}[H]
    \centering
    \includegraphics[width=0.72\linewidth]{cap02/osi_nodi_intermedi.png}
    \caption{I nodi intermedi implementano solo i primi tre strati; solo gli host implementano l'intera pila.}
\end{figure}

\begin{itemize}
    \item Solo i \textbf{sistemi terminali} (host) implementano l'intera pila, fino all'applicazione.
    \item I \textbf{nodi intermedi} (switch, router) implementano solo due o tre strati, i cosiddetti
          \emph{media layer}.
    \item Niente sopra il livello di rete: un router non ha alcun motivo di leggere il vostro messaggio.
\end{itemize}

\warning{Trasporto e rete: dove gira il software}{
    A livello di \textbf{trasporto} la comunicazione è tra \textbf{processi}, dal punto di vista logico è
    punto-punto, e il software è eseguito \textbf{sugli host}. A livello di \textbf{rete} la comunicazione
    è tra \textbf{host}: si risolve il problema dell'instradamento dei pacchetti attraverso la
    topologia della rete, e il software è eseguito prevalentemente \textbf{sui nodi della rete}. Un
    router non ha bisogno del trasporto perché non gli interessa a quale processo sia destinato il
    pacchetto: gli basta sapere a quale host deve arrivare.
}

% ══════════════════════════════════════════════════════════════
\section{Incapsulamento e decapsulamento}
% ══════════════════════════════════════════════════════════════

\begin{figure}[H]
    \centering
    \includegraphics[width=0.7\linewidth]{cap02/incapsulamento.png}
    \caption{Ogni strato del mittente aggiunge un header (e a volte un trailer); ogni strato del destinatario toglie il proprio.}
\end{figure}

\dfn{Incapsulamento e decapsulamento}{
    \begin{itemize}
      \item \textbf{Incapsulamento} (scendendo, \emph{top-down}): ogni strato del mittente aggiunge a ciò
            che riceve dallo strato superiore un proprio campo di controllo, sotto forma di un
            \textbf{header} (H) in testa e talvolta di un \textbf{trailer} (T) in coda.
      \item \textbf{Decapsulamento} (salendo, \emph{bottom-up}): ogni strato del destinatario rimuove e
            interpreta \textbf{esattamente} il campo aggiunto dal suo pari, e passa il resto verso l'alto.
    \end{itemize}
    Un'unità di dati, a qualunque strato, è quindi sempre \textbf{header + payload}, e il
    \textbf{payload} è l'unità dello strato superiore.
}

\begin{figure}[H]
\centering
\begin{tikzpicture}[every node/.style={font=\footnotesize\sffamily}]
  \node[field, fill=lapp, minimum width=4cm, minimum height=0.6cm] (m) at (0,0) {M};
  \node[anchor=west, lbl, text width=4.6cm, align=left] at (2.3,0) {messaggio \textbf{ricevuto} dal livello $k+1$, che deve essere inviato};
  \draw[-{Stealth}, very thick, netorange] (-2.4,-0.1) to[bend right=40] node[left, lbl] {livello $k$} (-2.4,-1.3);
  \node[field, fill=lapp, minimum width=4cm, minimum height=0.6cm] (m2) at (0,-1.4) {M};
  \node[field, fill=netorange!40, minimum width=0.8cm, minimum height=0.6cm, left=0pt of m2] {$H_k$};
  \node[anchor=west, lbl, text width=4.6cm, align=left] at (2.3,-1.4) {il livello $k$ aggiunge in testa il \textbf{proprio} header e passa tutto al livello $k-1$};
\end{tikzpicture}
\caption{Il passo elementare dell'incapsulamento, come alla lavagna: $M$ arriva dall'alto e diventa il payload di $H_k + M$.}
\end{figure}

\begin{figure}[H]
\centering
\begin{tikzpicture}[every node/.style={font=\scriptsize\sffamily}]
  % mittente
  \node[font=\small\bfseries\sffamily] at (2.2,1.6) {Mittente};
  \node[field, fill=lapp, minimum width=2cm] (m1) at (2.6,0) {M};
  \node[field, fill=ltra, minimum width=0.6cm, left=0pt of m1] (ht) {$H_t$};
  \node[field, fill=lapp, minimum width=2cm] (m2) at (2.6,-0.9) {M};
  \node[field, fill=ltra, minimum width=0.6cm, left=0pt of m2] (ht2) {$H_t$};
  \node[field, fill=lnet, minimum width=0.6cm, left=0pt of ht2] (hn2) {$H_n$};
  \node[field, fill=lapp, minimum width=2cm] (m3) at (2.6,-1.8) {M};
  \node[field, fill=ltra, minimum width=0.6cm, left=0pt of m3] (ht3) {$H_t$};
  \node[field, fill=lnet, minimum width=0.6cm, left=0pt of ht3] (hn3) {$H_n$};
  \node[field, fill=llink, minimum width=0.6cm, left=0pt of hn3] (hl3) {$H_l$};
  \node[field, fill=llink, minimum width=0.6cm, right=0pt of m3] (tl3) {$T_l$};
  \node[field, fill=lapp, minimum width=2cm, draw=none] (m0) at (2.6,0.9) {};
  % labels levels
  \foreach \y/\t/\u in {0.9/Applicazione/messaggio, 0/Trasporto/segmento, -0.9/Rete/datagramma, -1.8/Collegamento/frame, -2.7/Fisico/bit} {
     \node[anchor=east] at (-0.8,\y) {\t};
     \node[anchor=west, text=primary] at (4.4,\y) {\u};
  }
  \node[field, fill=lapp, minimum width=2cm] at (2.6,0.9) {M};
  \node at (2.2,-2.7) {0110100101110\dots};
  \draw[-{Stealth}, thick, netred] (-0.3,1.1) -- (-0.3,-2.9);
  % destinatario
  \begin{scope}[xshift=8.5cm]
    \node[font=\small\bfseries\sffamily] at (2.2,1.6) {Destinatario};
    \node[field, fill=lapp, minimum width=2cm] at (2.6,0.9) {M};
    \node[field, fill=lapp, minimum width=2cm] (d1) at (2.6,0) {M};
    \node[field, fill=ltra, minimum width=0.6cm, left=0pt of d1] {$H_t$};
    \node[field, fill=lapp, minimum width=2cm] (d2) at (2.6,-0.9) {M};
    \node[field, fill=ltra, minimum width=0.6cm, left=0pt of d2] (e2) {$H_t$};
    \node[field, fill=lnet, minimum width=0.6cm, left=0pt of e2] {$H_n$};
    \node[field, fill=lapp, minimum width=2cm] (d3) at (2.6,-1.8) {M};
    \node[field, fill=ltra, minimum width=0.6cm, left=0pt of d3] (e3) {$H_t$};
    \node[field, fill=lnet, minimum width=0.6cm, left=0pt of e3] (f3) {$H_n$};
    \node[field, fill=llink, minimum width=0.6cm, left=0pt of f3] {$H_l$};
    \node[field, fill=llink, minimum width=0.6cm, right=0pt of d3] {$T_l$};
    \node at (2.2,-2.7) {0110100101110\dots};
    \draw[-{Stealth}, thick, netgreen] (4.6,-2.9) -- (4.6,1.1);
  \end{scope}
  \draw[-{Stealth}, very thick, black!60, decorate, decoration={snake, amplitude=1.2pt, segment length=6pt, post length=2mm}] (5.4,-2.7) -- node[below=3pt] {mezzo fisico} (9.5,-2.7);
\end{tikzpicture}
\caption{Incapsulamento nella pila a cinque strati: ogni strato tratta tutto ciò che riceve dall'alto come un blocco unico a cui aggiungere il proprio header.}
\end{figure}

\subsection{Sotto il proprio strato, il payload è opaco}

Quando il messaggio di livello 7 scende, tutto ciò che viene dall'alto (header e contenuto insieme)
diventa il \textbf{payload} del livello 6.

\begin{itemize}
    \item Il livello 6 aggiunge il proprio header e passa tutto al livello 5, che fa lo stesso, e così via.
    \item Il livello 6 \textbf{non sa} di trasportare HTML, e non deve interessargli: per lui il payload è
          una \textbf{sequenza di byte opaca}.
    \item È proprio questo che rende gli strati \textbf{indipendenti}: uno strato che ispezionasse il
          proprio payload smetterebbe di funzionare appena cambia lo strato superiore.
\end{itemize}

Il destinatario ripercorre la sequenza al contrario, e ogni strato toglie soltanto il proprio header.

\info{Lo tocchiamo con mano}{
    Quando nel capitolo sui socket scriveremo a mano una richiesta HTTP, vedremo che un messaggio a
    livello di applicazione viene semplicemente inviato come \textbf{payload} di un messaggio di
    trasporto, il cui contenuto è del tutto irrilevante (``trasparente'') per il livello di trasporto.
}

% ══════════════════════════════════════════════════════════════
\section{I sette strati, uno per uno}
% ══════════════════════════════════════════════════════════════

\subsection{1. Livello fisico (Physical)}

\begin{figure}[H]
\centering
\begin{minipage}{0.48\linewidth}\centering
  \includegraphics[width=\linewidth]{cap02/osi_fisico.png}
\end{minipage}\hfill
\begin{minipage}{0.48\linewidth}
È responsabile di \textbf{spostare i singoli bit} da un nodo al successivo lungo il mezzo. Si occupa di:
\begin{itemize}
  \item caratteristiche fisiche di interfacce e mezzi (connettori, cavi, segnali elettrici);
  \item configurazione della linea (punto-punto o multipunto) e topologia fisica;
  \item modalità di trasmissione (simplex, half-duplex, full-duplex);
  \item rappresentazione dei bit, velocità (data rate) e sincronizzazione dei bit tra i due estremi.
\end{itemize}
\end{minipage}
\end{figure}

\begin{figure}[H]
\centering
\begin{tikzpicture}[every node/.style={font=\footnotesize\sffamily}]
  \node (t1) at (0,0.9) {\icon[0.8cm]{pc}}; \node[above=-1pt of t1, text=netblue, font=\small\bfseries\sffamily] {T1};
  \node (t2) at (10,0.9) {\icon[0.8cm]{pc}}; \node[above=-1pt of t2, text=netblue, font=\small\bfseries\sffamily] {T2};
  \node[font=\ttfamily\footnotesize, text=netblue] (b1) at (0,0) {0001100110000110};
  \node[font=\ttfamily\footnotesize, text=netblue] (b2) at (10,0) {0001100110000110};
  \node[draw, fill=lphy, minimum width=1.2cm, minimum height=0.5cm] (p1) at (1.2,-1) {PHYS};
  \node[draw, fill=lphy, minimum width=1.2cm, minimum height=0.5cm] (p2) at (8.8,-1) {PHYS};
  \node[cylinder, shape border rotate=0, draw, thick, fill=black!5, aspect=0.25, minimum height=4.8cm, minimum width=0.55cm] (c) at (5,-1) {};
  \node[below=2pt of c, lbl] {mezzo fisico: cavo, fibra, onde radio};
  \draw[-{Stealth}, thick, netorange] (b1.south) to[bend right=20] (p1.north);
  \draw[-{Stealth}, thick, netorange] (p1) -- (c.west);
  \draw[-{Stealth}, thick, netorange] (c.east) -- (p2);
  \draw[-{Stealth}, thick, netorange] (p2.north) to[bend right=20] (b2.south);
  \foreach \x in {3.3,3.9,4.5,5.1,5.7,6.3} \draw[netorange!80!black, thick, decorate, decoration={snake, amplitude=1.2pt, segment length=4pt}] (\x,-1) -- ++(0.4,0);
\end{tikzpicture}
\caption{Il livello fisico: la sequenza di bit di T1 diventa un segnale sul mezzo e ricompare identica (se tutto va bene) in T2.}
\end{figure}

\subsection{2. Livello di collegamento (Data Link)}

\begin{figure}[H]
\centering
\begin{minipage}{0.48\linewidth}\centering
  \includegraphics[width=\linewidth]{cap02/osi_collegamento.png}
\end{minipage}\hfill
\begin{minipage}{0.48\linewidth}
È responsabile di \textbf{spostare i frame da un nodo al successivo}: un solo salto (\emph{hop}),
non l'intero cammino. Si occupa di:
\begin{itemize}
  \item organizzare i dati in \textbf{frame}, blocchi di dimensione limitata (tipicamente qualche
        centinaio di byte), con header $H_2$ e trailer $T_2$;
  \item \textbf{controllo di flusso} e \textbf{controllo degli errori}, tramite la \emph{frame check
        sequence} (FCS) nel trailer;
  \item \textbf{controllo dell'accesso} al mezzo quando è condiviso.
\end{itemize}
\end{minipage}
\end{figure}

\dfn{Sottostrato MAC}{
    Il livello di collegamento contiene il \textbf{sottostrato MAC} (\emph{Medium Access Control}), che si occupa
    della \textbf{gestione dell'accesso a un unico mezzo di trasmissione} da parte di più stazioni \textbf{in
    competizione} tra loro. Ethernet ne è l'esempio più noto.
}

\subsection{3. Livello di rete (Network)}

\begin{figure}[H]
\centering
\begin{minipage}{0.48\linewidth}\centering
  \includegraphics[width=\linewidth]{cap02/osi_rete.png}
\end{minipage}\hfill
\begin{minipage}{0.48\linewidth}
È responsabile di \textbf{consegnare i pacchetti dall'host sorgente all'host destinazione}, lungo
l'intero cammino. Si occupa di:
\begin{itemize}
  \item attraversare la \textbf{topologia}: la rete è un grafo di nodi con più di una strada possibile;
  \item \textbf{indirizzamento logico} (gli indirizzi IP);
  \item \textbf{routing}: decidere, nodo per nodo, da quale collegamento far uscire il pacchetto
        (tabelle di instradamento).
\end{itemize}
\end{minipage}
\end{figure}

\begin{figure}[H]
\centering
\begin{tikzpicture}[every node/.style={font=\footnotesize\sffamily}]
  \begin{scope}
    \node[font=\small\bfseries\sffamily, fill=llink!70, rounded corners] at (1.6,1.5) {ISO-2: collegamento};
    \node[draw=netred, thick, minimum width=0.9cm, minimum height=0.55cm] (a) at (0,0) {};
    \node[draw=netred, thick, minimum width=0.9cm, minimum height=0.55cm] (b) at (3.2,0) {};
    \draw[{Stealth}-{Stealth}, dashed, thick, netred] (a) -- (b);
    \node[field, fill=llink, minimum width=0.9cm] at (1.6,0.45) {frame};
    \node[lbl, text width=4.2cm] at (1.6,-1) {ci sono \textbf{solo due nodi} che comunicano tra loro; viaggiano frame di qualche centinaio di byte};
  \end{scope}
  \begin{scope}[xshift=7.2cm]
    \node[font=\small\bfseries\sffamily, fill=lnet!70, rounded corners] at (1.8,1.5) {ISO-3: rete};
    \node[gnode] (s) at (0,0.3) {}; \node[left=1pt of s] {S};
    \node[gnode] (r1) at (1.2,0.9) {}; \node[gnode] (r2) at (1.3,-0.5) {};
    \node[gnode] (r3) at (2.5,0.4) {}; \node[gnode] (r4) at (2.7,-0.9) {};
    \node[gnode] (d) at (3.8,0) {}; \node[right=1pt of d] {D};
    \foreach \x/\y in {s/r1, s/r2, r1/r3, r2/r3, r2/r4, r3/d, r4/d, r1/r2} \draw[black!40, thick] (\x) -- (\y);
    \draw[pathhl, opacity=0.6] (s) -- (r2) -- (r3) -- (d);
    \node[lbl, text width=5.2cm, anchor=north] at (1.9,-1.25) {chi trasmette attraverso una \textbf{topologia complessa}? Il livello di rete, usando i servizi del collegamento su ogni salto};
  \end{scope}
\end{tikzpicture}
\caption{Il collegamento lavora su un solo salto tra due nodi; la rete compone i salti in un cammino dalla sorgente alla destinazione.}
\end{figure}

\subsection{4. Livello di trasporto (Transport)}

\begin{figure}[H]
\centering
\begin{minipage}{0.48\linewidth}\centering
  \includegraphics[width=\linewidth]{cap02/osi_trasporto.png}
\end{minipage}\hfill
\begin{minipage}{0.48\linewidth}
È responsabile di \textbf{consegnare un messaggio da un estremo all'altro} e al \textbf{programma
giusto} sull'host di destinazione. Si occupa di:
\begin{itemize}
  \item \textbf{indirizzamento tramite porte}: il livello di rete trova l'\emph{host}, il trasporto
        trova il \emph{processo};
  \item controllo della connessione (con o senza connessione);
  \item \textbf{segmentazione e riassemblaggio} dei messaggi in segmenti;
  \item controllo di flusso e degli errori, da un estremo all'altro.
\end{itemize}
\end{minipage}
\end{figure}

\subsection{5. Livello di sessione (Session)}

\begin{figure}[H]
\centering
\begin{minipage}{0.48\linewidth}\centering
  \includegraphics[width=\linewidth]{cap02/osi_sessione.png}
\end{minipage}\hfill
\begin{minipage}{0.48\linewidth}
È responsabile del \textbf{controllo del dialogo}: stabilisce, mantiene e sincronizza l'interazione
tra i due sistemi.
\begin{itemize}
  \item \textbf{chi parla, quando e per quanto};
  \item \textbf{punti di sincronizzazione} (\emph{checkpoint}, i \texttt{syn} della figura): uno scambio
        interrotto riprende dall'ultimo checkpoint invece che dall'inizio.
\end{itemize}
\end{minipage}
\end{figure}

\subsection{6. Livello di presentazione (Presentation)}

\begin{figure}[H]
\centering
\begin{minipage}{0.48\linewidth}\centering
  \includegraphics[width=\linewidth]{cap02/osi_presentazione.png}
\end{minipage}\hfill
\begin{minipage}{0.48\linewidth}
È responsabile della \textbf{rappresentazione dei dati}, perché ciò che scrive una parte sia ciò che
legge l'altra:
\begin{itemize}
  \item \textbf{traduzione} tra codifiche diverse (il caso storico è ASCII contro EBCDIC);
  \item \textbf{cifratura} e decifratura;
  \item \textbf{compressione}.
\end{itemize}
\end{minipage}
\end{figure}

\subsection{7. Livello applicativo (Application)}

\begin{figure}[H]
\centering
\begin{minipage}{0.48\linewidth}\centering
  \includegraphics[width=\linewidth]{cap02/osi_applicazione.png}
\end{minipage}\hfill
\begin{minipage}{0.48\linewidth}
È responsabile di \textbf{fornire servizi all'utente} (una persona o un programma), ed è dove vivono i
protocolli che già conoscete per nome:
\begin{itemize}
  \item login remoto (terminale virtuale di rete);
  \item trasferimento e accesso ai file;
  \item posta elettronica;
  \item accesso al World Wide Web.
\end{itemize}
È lo strato che apriremo per primo, perché è quello che si vede.
\end{minipage}
\end{figure}

% ══════════════════════════════════════════════════════════════
\section{Perché OSI non ha vinto}
% ══════════════════════════════════════════════════════════════

Il modello OSI è lo standard \emph{de iure}, ma i suoi protocolli non hanno avuto successo. Le
ragioni sono quattro:

\begin{itemize}
    \item \textbf{Cattivo tempismo} (\emph{bad timing}): quando comparvero i protocolli OSI, TCP/IP era
          già diffuso nelle università di ricerca, e nessun produttore voleva sostenere un secondo stack.
    \item \textbf{Cattiva progettazione} (\emph{bad design}): sia il modello sia i protocolli avevano
          difetti. Sette strati erano un numero ``politico'': sessione e presentazione sono quasi vuoti,
          collegamento e rete sono sovraccarichi.
    \item \textbf{Cattive implementazioni} (\emph{bad implementations}): le prime erano enormi, goffe e
          lente, e la cattiva reputazione è rimasta anche dopo che i prodotti erano migliorati.
    \item \textbf{Cattiva politica} (\emph{bad politics}): era percepito come una creatura dei ministeri
          delle telecomunicazioni europei, della Comunità Europea e poi del governo statunitense,
          mentre TCP/IP arrivava gratis insieme a Berkeley UNIX.
\end{itemize}

Il \textbf{modello} OSI è sopravvissuto ed è ancora il modo in cui si parla di reti; i suoi
\textbf{protocolli} no.

\subsection{L'apocalisse dei due elefanti}

\begin{figure}[H]
    \centering
    \includegraphics[width=0.5\linewidth]{cap02/due_elefanti.png}
    \caption{L'attività su un argomento: prima la ricerca, poi il miliardo di dollari di investimenti.}
\end{figure}

Quando nasce un nuovo argomento, l'attività ha due ``gobbe'' (i due elefanti): prima un picco di
\textbf{ricerca}, poi, anni dopo, un picco di \textbf{investimenti industriali}. Tra le due c'è una conca.

\begin{itemize}
    \item Uno standard scritto \textbf{troppo presto} congela un argomento che nessuno ha ancora capito.
    \item Uno standard scritto \textbf{troppo tardi} viene ignorato da chi ha già investito in altro.
\end{itemize}

Per evitare il cattivo tempismo, gli standard vanno scritti \textbf{nella conca tra le due gobbe}.
OSI non ci riuscì: arrivò quando TCP/IP aveva già occupato il secondo elefante.

% ══════════════════════════════════════════════════════════════
\section{Cosa non va nel modello TCP/IP}
% ══════════════════════════════════════════════════════════════

Per correttezza, anche il \emph{modello} TCP/IP ha i suoi difetti:

\begin{itemize}
    \item \textbf{non distingue chiaramente} i concetti di servizio, interfaccia e protocollo, cioè
          proprio ciò che OSI aveva fatto bene;
    \item \textbf{non è generale}: non si riesce a usarlo per descrivere altre pile di protocolli
          (provate a descrivere Bluetooth con il modello TCP/IP);
    \item il suo \textbf{livello di collegamento non è un vero strato}, ma un'interfaccia tra l'host e il
          collegamento trasmissivo;
    \item \textbf{non distingue} tra livello fisico e livello di collegamento, che sono due compiti
          completamente diversi.
\end{itemize}

I suoi \textbf{protocolli}, però, erano progettati con cura, ben implementati e distribuiti
\textbf{gratuitamente}.

% ══════════════════════════════════════════════════════════════
\section{De iure e de facto}
% ══════════════════════════════════════════════════════════════

\begin{figure}[H]
\centering
\begin{minipage}{0.45\linewidth}\centering
  \includegraphics[width=\linewidth]{cap02/de_iure_de_facto.png}
\end{minipage}\hfill
\begin{minipage}{0.5\linewidth}
\begin{itemize}
  \item \textbf{De iure}: ISO/OSI è il modello standard per le reti di calcolatori: completo,
        dettagliato e complesso.
  \item \textbf{De facto}: i protocolli che girano davvero su Internet e sulle reti locali sono
        \textbf{TCP}, \textbf{UDP} e \textbf{IP} (la \emph{Internet Protocol Suite}):
        \begin{itemize}
          \item \textbf{TCP}: \emph{Transmission Control Protocol};
          \item \textbf{UDP}: \emph{User Datagram Protocol};
          \item \textbf{IP}: \emph{Internet Protocol}.
        \end{itemize}
\end{itemize}
In pratica non sono rivali: teniamo il \textbf{modello} per ragionare e i \textbf{protocolli} per lavorare.
\end{minipage}
\caption{Corrispondenza tra la pila ISO/OSI e la pila TCP/IP.}
\end{figure}

\info{IP è il collante}{
    Il protocollo IP è il ``collante'' di tutta la rete di reti che chiamiamo Internet. È l'unico punto
    in cui tutti sono obbligati a mettersi d'accordo: sopra IP ci possono essere TCP o UDP, sotto IP
    Ethernet, Wi-Fi o qualunque altra cosa. È la ``vita stretta'' della clessidra vista nel capitolo 1.
}

% ══════════════════════════════════════════════════════════════
\section{I cinque strati usati nel resto del corso}
% ══════════════════════════════════════════════════════════════

Da qui in avanti il corso conta \textbf{cinque strati}, non sette: è il modello OSI senza sessione e
presentazione, i cui pochi compiti reali vengono svolti direttamente dalle applicazioni. I nomi sono
quelli del Kurose.

\begin{figure}[H]
\centering
\begin{tikzpicture}[every node/.style={font=\small\sffamily}]
  \node[lay app, minimum width=3.4cm, minimum height=0.8cm] (a) at (0,2.4) {Applicazione};
  \node[lay tra, minimum width=3.4cm, minimum height=0.8cm] (t) at (0,1.6) {Trasporto};
  \node[lay net, minimum width=3.4cm, minimum height=0.8cm] (n) at (0,0.8) {Rete};
  \node[lay link, minimum width=3.4cm, minimum height=0.8cm] (l) at (0,0) {Collegamento};
  \node[lay phy, minimum width=3.4cm, minimum height=0.8cm] (p) at (0,-0.8) {Fisico};
  \foreach \n/\osi/\u/\ex in {a/{OSI 7, 6, 5}/messaggio/{HTTP, SMTP, DNS}, t/{OSI 4}/segmento/{TCP, UDP}, n/{OSI 3}/datagramma/{IP}, l/{OSI 2}/frame/{Ethernet, Wi-Fi}, p/{OSI 1}/bit/{rame, fibra, radio}} {
     \node[anchor=east, lbl] at ([xshift=-0.3cm]\n.west) {\osi};
     \node[anchor=west, text=primary!85!black, font=\small\bfseries\sffamily] at ([xshift=0.3cm]\n.east) {\u};
     \node[anchor=west, lbl, text=black!60] at ([xshift=2.4cm]\n.east) {\ex};
  }
  \node[font=\footnotesize\bfseries\sffamily] at (2.7,3.2) {unità di dati};
  \node[font=\footnotesize\bfseries\sffamily] at (4.7,3.2) {esempi};
\end{tikzpicture}
\caption{La pila a cinque strati: corrispondenza con OSI, nome dell'unità di dati ed esempi di protocolli.}
\end{figure}

\begin{table}[H]
\centering
\begin{tabular}{lcl}
\toprule
\textbf{Strato} & \textbf{Corrisponde a OSI} & \textbf{Unità di dati} \\
\midrule
Applicazione & 7, 6 e 5 & messaggio \\
Trasporto & 4 & segmento \\
Rete & 3 & datagramma \\
Collegamento & 2 & frame \\
Fisico & 1 & bit \\
\bottomrule
\end{tabular}
\caption{I cinque strati del corso.}
\end{table}

\warning{Numeri OSI e nomi del corso}{
    I numeri OSI (``livello 3'', ``livello 2'') restano in uso perché li usano tutti nel settore, e li
    useremo anche noi quando serve. Ma la pila che studiamo, dall'alto verso il basso, è quella a
    cinque strati.
}

% ══════════════════════════════════════════════════════════════
\section{La stratificazione oltre Internet}
% ══════════════════════════════════════════════════════════════

Esistono altre pile di protocolli, scelte a seconda del contesto: \textbf{RS-232} per una linea
seriale, \textbf{CAN} all'interno di un veicolo o in automazione industriale. Sono molto più semplici,
e sono comunque stratificate.

\begin{figure}[H]
    \centering
    \includegraphics[width=0.75\linewidth]{cap02/altri_modelli.png}
    \caption{Modelli a strati in altri contesti (standard industriali, automazione, IoT).}
\end{figure}

\begin{table}[H]
\centering
\begin{tabular}{cccc}
\toprule
\textbf{Strato OSI} & \textbf{TCP/IP o UDP/IP} & \textbf{RS-232} & \textbf{CAN} \\
\midrule
7--5 & Applicazione & \multirow{3}{*}{Applicazione} & \multirow{3}{*}{\makecell{CANopen, CANopen FD,\\ J1939, \dots}} \\
4 & TCP/UDP & & \\
3 & IP & & \\
\midrule
2 & \multirow{2}{*}{Link + Fisico} & UART & CAN 2.0 A/B \\
1 & & Rx, Tx, GND & CAN-H, CAN-L \\
\bottomrule
\end{tabular}
\caption{In RS-232 e CAN i livelli di collegamento e fisico sono distinti, e tutto il resto è ``applicazione''.}
\end{table}

% ══════════════════════════════════════════════════════════════
\section{Riepilogo}
% ══════════════════════════════════════════════════════════════

\riepilogo{
\begin{itemize}
    \item Una richiesta attraversa troppi pezzi per essere progettata tutta insieme: il sistema si divide
          in \textbf{strati} (\emph{divide et impera}).
    \item Uno strato offre un \textbf{servizio} verso l'alto attraverso un'\textbf{interfaccia}, e usa un
          \textbf{protocollo} in orizzontale con il suo pari. Le due cose sono indipendenti.
    \item La comunicazione tra pari è \textbf{logica}: i dati scendono una pila e risalgono l'altra, e solo
          il livello fisico sposta davvero qualcosa.
    \item I servizi possono essere \textbf{con} o \textbf{senza connessione}, \textbf{affidabili} o no:
          TCP darà un flusso di byte affidabile, UDP un datagram inaffidabile.
    \item L'\textbf{incapsulamento} rende possibile tutto questo: ogni strato tratta ciò che trasporta come
          byte opachi.
    \item OSI ci ha dato sette strati con cui ragionare, TCP/IP i protocolli; nel corso usiamo
          \textbf{cinque strati} con i nomi del Kurose: messaggio, segmento, datagramma, frame, bit.
\end{itemize}
}

Nella prossima parte apriamo la cima della pila, dove stanno le applicazioni.
